\section{数据理解与总体分析}
\subsection{输入结构与审计}
输入同时提供节点和边的 JSON、CSV 表示。节点编号用于识别对象，不代表最优执行次序。操作节点的缓冲区列表同时包含输入与输出，不能仅依据列表顺序推断读写属性。本文直接使用题面给出的依赖边约束生产与消费先后。

六图规模如下。读取时逐字段比对两种表示，检查节点编号、重复边、缓存类型、大小及申请释放配对。全部输入通过检查，没有对缺失值进行填补，也没有删除任何节点或依赖。
\begin{table}[H]\centering\caption{输入计算图规模}\small
\begin{tabular}{lrr}\toprule 计算图 & 节点数 & 依赖边数\\\midrule
Conv\_Case0 & 2580 & 3869\\
Conv\_Case1 & 36086 & 85653\\
FlashAttention\_Case0 & 1716 & 2712\\
FlashAttention\_Case1 & 6952 & 11184\\
Matmul\_Case0 & 4160 & 7104\\
Matmul\_Case1 & 30976 & 55040\\\bottomrule
\end{tabular}\end{table}

\subsection{目标的耦合关系}
缓存驻留峰值低，并不意味着每种缓存都能分配成功。一方面，不同类型不能互相借用空间；另一方面，分散的空闲区间无法容纳较大的连续申请。与此相反，过度追求小峰值可能使原本可重叠的计算串行化。因此首先通过生命周期收缩建立基准，再分别围绕搬运量和时间评价候选。

采用“拓扑顺序—连续分配—依赖重建—流水评价”的计算链。求解程序只产生候选，随后由不调用求解模块的重放程序从文本附件读回结果，重新检查空间和时间约束。两个程序共享输入字段定义，但不共享区间分配逻辑，减少同一实现错误同时影响生成和检查的可能性。
